% -----------------------------------------------------------------------------
% This file is automatically generated.
% Do not edit manually.
% Source: notebooks/support/population-estimation/support.Rmd
% -----------------------------------------------------------------------------


\noindent We saw in Section 3.2.3 that the finiteness of the population from which we sample has a noticeable effect when the sample size is greater than 10\% of the population size. In the example of the mean monthly payroll ($n = 50$, $N = 2,222$), it can be expected that this is the case if we want the precision around the total monthly payroll to be less than 300,000.

\noindent The confidence interval incorporating the finite-population correction factor is given by Equation 3.10:

\par\addvspace{\topsep}
\begingroup
\fvset{listparameters={%
  \setlength{\topsep}{0pt}%
  \setlength{\partopsep}{0pt}%
  \setlength{\parsep}{0pt}%
  \setlength{\itemsep}{0pt}%
}}
\begin{Verbatim}[breaklines=true,formatcom=\color{ada_blue}]
E <- 300000
tval <- qt(0.975, df = (n - 1))
(gamma <- E^2 / (N * tval^2 * sd(sample1$gross)^2))
\end{Verbatim}
\begin{Verbatim}[breaklines=true,formatcom=\color{ada_light_blue}]
[1] 4.242586
\end{Verbatim}
\begin{Verbatim}[breaklines=true,formatcom=\color{ada_blue}]
(N / (1 + gamma))
\end{Verbatim}
\begin{Verbatim}[breaklines=true,formatcom=\color{ada_light_blue}]
[1] 423.8366
\end{Verbatim}
\endgroup
\par\addvspace{\topsep}

\noindent The total sample size required to obtain an estimate of the mean monthly payroll with a precision of 300,000 at 95\% confidence is 424. This is more than 10\% of the population size; therefore, we were right to use the finite-population correction factor. If we had not anticipated that the required sample size would have been unnecessarily large, not taking the finite-population correction factor into account would have resulted in a sample size of 524.
